% Lösungen Ableitung und Änderungsrate · Mittlere Änderungsrate · Prüfungs-Fokus Abitur GK (zusammenbau v0.6)
\einheitenkopf{Einheit 1 · Mittlere Änderungsrate und Sekante}

\erg{1}{a) Zeitraum – Differenzenquotient, Sekante: „in den ersten fünf Stunden“ nennt ein Intervall \quad b) $\frac{300 - 120}{3} = 60$ Besucher je Stunde \quad c) $A$: $\frac{A(10) - A(0)}{10} \approx -4{,}42$; $B$: $\frac{60 - 90}{10} = -3$; der Betrag ist bei $A$ größer}
\erg{2}{a) $a(1) = 100$, $a(3) = 900$; $\frac{900 - 100}{2} = 400$ Bewertungen je Stunde \quad b) $h(5) \approx 63{,}36$, $h(0) = 40$; $\frac{h(5) - h(0)}{5} \approx 4{,}67$ Prozentpunkte je Tag \quad c) $k(16) = 32{,}6$, $k(10) = 5$; $\frac{k(16) - k(10)}{6} = 4{,}6$ mmol/l je km/h \quad d) I: $2 - 1{,}1 = 0{,}9$ Liter – so viel Luft wird beim Einatmen zwischen $1{,}5$ und $3$ Sekunden aufgenommen; II: $0{,}9 : 1{,}5 = 0{,}6$ Liter je Sekunde – die mittlere Einatemrate in diesem Zeitraum \quad e) $f(1) \approx 2{,}011$, $f(5) \approx 0{,}947$; $m = \frac{f(5) - f(1)}{4} \approx -0{,}266$; $y = -0{,}266 \cdot (x - 1) + 2{,}011 \approx -0{,}266x + 2{,}277$ \quad f) $f$: $\frac{f(5) - f(0)}{5} \approx -0{,}211$; $h$: $\frac{0{,}6 - 1{,}8}{5} = -0{,}24$; der Betrag der mittleren Steigung ist bei $h$ größer ($0{,}24 > 0{,}211$)}
\erg{3}{a) Länge $4 \cdot 3{,}2 + 2\pi \cdot 0{,}5 \approx 15{,}94$ m; Zeit $\frac{15{,}94}{4} \approx 3{,}99$ min \quad b) $\frac{k(c) - k(0)}{c} = 0{,}1c^2 - 1{,}5c = -5$, also $c^2 - 15c + 50 = 0$ und $c = 5$ oder $c = 10$; kleinere Lösung $c = 5$: Die linke Seite ist die mittlere Abbaurate über $[0; c]$; in den ersten fünf Jahren nimmt der Bestand im Mittel um $5$ Mio. Tonnen je Jahr ab. \quad c) Falsch: $K(2) - K(1) = 10$, aber $K(3) - K(2) = 4$ – die Kosten des zusätzlichen Kubikmeters nehmen zunächst ab, obwohl $K$ steigt. \quad d) Falsch: $K(2) - K(1) = 15{,}5$, aber $K(3) - K(2) = 9{,}5$ – die Kosten des zusätzlichen Kubikmeters nehmen zunächst ab, obwohl $K$ steigt. \quad e) $g(x + 1) - g(x) = -500\mathrm{e}^{-0{,}4x} \cdot \mathrm{e}^{-0{,}4} + 500\mathrm{e}^{-0{,}4x} = 500(1 - \mathrm{e}^{-0{,}4})\mathrm{e}^{-0{,}4x}$ (Intervalllänge $1$); $500(1 - \mathrm{e}^{-0{,}4}) \approx 164{,}8$; $\mathrm{e}^{-0{,}4x} = \frac{80}{164{,}8}$ liefert $x \approx 1{,}81$ h, also nach etwa 1 h 48 min \quad f) $g(x + 1) - g(x) = -400\mathrm{e}^{-0{,}3x} \cdot \mathrm{e}^{-0{,}3} + 400\mathrm{e}^{-0{,}3x} = 400(1 - \mathrm{e}^{-0{,}3})\mathrm{e}^{-0{,}3x}$ (Intervalllänge $1$); $400(1 - \mathrm{e}^{-0{,}3}) \approx 103{,}7$; $\mathrm{e}^{-0{,}3x} = \frac{60}{103{,}7}$ liefert $x \approx 1{,}82$ h, also nach etwa 1 h 49 min}
